% =========================================================
% Proof Stub: Ratio Theorem for Sequences
% Source: volume-iii/analysis/sequences/notes/asymptotics/notes-growth-and-asymptotics.tex
% =========================================================

% *** HARDER PROOF — attempt after foundational results settled ***

\clearpage
\phantomsection
\hypertarget{proof-RA-ABB-T-ratio-thm-seq}{}

\subsubsection[Ratio Theorem for Sequences]{Proof --- Ratio Theorem for Sequences}

\bigskip

\noindent
\textbf{Source.} See \texttt{notes-growth-and-asymptotics.tex}.

\vspace{0.75em}

\noindent
\textbf{Goal.} If $a_n > 0$ and $a_{n+1}/a_n \to L$, then $a_n^{1/n} \to L$.

\vspace{0.75em}

\noindent
\textbf{Logical form.}
\[
  \frac{a_{n+1}}{a_n} \to L \Rightarrow a_n^{1/n} \to L.
\]

\vspace{0.75em}

\noindent
\textbf{Proof strategy.} Apply the geometric-arithmetic mean estimate: express $a_n^{1/n}$ using telescoping products, then use the convergence of the ratio to bound $a_n^{1/n}$ between $(L-\varepsilon)$ and $(L+\varepsilon)$ for large $n$.

\vspace{0.75em}

\noindent
\textbf{Proof.}
\begin{proof}
\Claim If $a_n > 0$ and $a_{n+1}/a_n \to L$, then $a_n^{1/n} \to L$.

\Given 

\Goal 

\Strategy Apply the geometric-arithmetic mean estimate: express $a_n^{1/n}$ using telescoping products, then use the convergence of the ratio to bound $a_n^{1/n}$ between $(L-\varepsilon)$ and $(L+\varepsilon)$ for large $n$.

\medskip

% --- write a_n = a_1 · ∏_{k=1}^{n-1} (a_{k+1}/a_k) ---
% --- since a_{k+1}/a_k → L: for k ≥ N, L-ε < a_{k+1}/a_k < L+ε ---
% --- split product at N; bound tail by (L±ε)^{n-N} ---
% --- take n-th root; (L+ε)^{(n-N)/n} → L+ε as n → ∞ ---
% --- conclude limsup a_n^{1/n} ≤ L+ε and liminf ≥ L-ε; let ε→0 ---

\end{proof}

\begin{remark*}[Proof shape]
Control the telescoping product of ratios; taking $n$-th roots then transfers the ratio limit to a root limit.
\end{remark*}

\begin{remark*}[Dependencies]
\begin{itemize}
  \item Algebra of limits.
  \item Geometric-arithmetic mean bound on products.
  \item For bounded sequences: controlling limsup and liminf is sufficient to find the limit.
\end{itemize}
\end{remark*}
